\chapter{METHODOLOGY}\label{ch:method}

\begin{doublespacing}

This chapter presents the four algorithmic models and the Hungarian
baselines, then derives upper bounds on their expected finder
competitive ratio. Table~\ref{tab:models} summarizes the models along
the two operational axes of energy and communication.

\begin{table}[ht]
\centering
\caption{Summary of the four algorithmic models and their operational characteristics.}\label{tab:models}
\vspace{2pt}
\begin{tabular}{@{}lcccc@{}}
\toprule
& \textbf{Energy} & \textbf{Communication} & \textbf{Movement} & \textbf{Coordination} \\
\midrule
M1 & unbounded & none & Base $\to$ node $\to$ base & Random \\
M2 & unbounded & full & Node $\to$ node & Auction \\
M3 & $E$ & at base & Base $\to$ node $\to$ base & Auction \\
M4 & $E$ & at base & Base $\to$ tour $\to$ base & Auction \\
\bottomrule
\end{tabular}
\end{table}

The four models are designed to answer the question ``how much does
each additional capability help?'' Moving from M1 to M2 isolates the
value of communication and coordination. Moving from M2 to M3 isolates
the cost of a finite energy budget. Moving from M3 to M4 isolates the
value of chaining several sites into a single sortie. The relationship
is shown graphically in Figure~\ref{fig:model-tree}.

\begin{figure}[ht]
  \centering
  \includegraphics[width=0.85\textwidth]{fig_model_tree.png}
  \caption{The four models ordered by what each adds. M4* (the bid function $p_i/d^2$ variant of M4) is the primary contribution of the thesis.}\label{fig:model-tree}
\end{figure}

\subsection{Instance Generator}

To reproduce any result in this thesis, a reader needs only the instance
generator and the model loop described below. The instance generator
creates a single search scenario: it places $n$ sites and $R$ robot bases
uniformly at random in the square $[0, L]^2$, draws prior probabilities
from a uniform distribution and normalizes them to sum to one, then
selects a target site by sampling from those priors. Two feasibility
filters are applied: the nearest-base distance must be at least
$d_{\min} = 1.0$ (to avoid degenerate near-zero FCR denominators), and
the target must be reachable within the energy budget ($2\,d_{\mathrm{opt}}
\leq E$). The following Python pseudocode implements this exactly:

\begin{verbatim}
def generate_instance(n, R, L, E, d_min, seed):
    rng = np.random.default_rng(seed)       # reproducible random state

    # Place n sites uniformly in [0,L]^2 and draw normalized priors
    sites = rng.uniform(0, L, size=(n, 2))
    raw   = rng.uniform(0.1, 1.0, size=n)
    p     = raw / raw.sum()                 # prior probabilities sum to 1

    # Place R robot bases uniformly in [0,L]^2
    bases = rng.uniform(0, L, size=(R, 2))

    # Rejection-sample a feasible target site
    for _ in range(10_000):
        tau   = rng.choice(n, p=p)          # sample target from prior
        d_opt = min(np.linalg.norm(bases[r] - sites[tau])
                    for r in range(R))      # nearest base distance
        if d_opt >= d_min and 2 * d_opt <= E:
            return sites, p, bases, tau     # feasible instance found
    raise ValueError("No feasible instance found for given parameters")
\end{verbatim}

Each call to \texttt{generate\_instance} with the same seed produces an
identical scenario, which is how all models are evaluated on the same
instances for a fair paired comparison. The 5{,}000-trial main comparison
uses seeds $42, 43, \ldots, 5041$; every result in Chapter~V can be
exactly reproduced by running \texttt{python experiments/exp8\_ci\_table.py}
from the repository root.

\section{Algorithmic Models}

\subsection{Model 1: Random Baseline}

\paragraph{The idea.} Model 1 is deliberately a bad algorithm, chosen
to measure how much a team loses when its members cannot coordinate.
Each robot, on its own, picks a site at random (weighted by how
likely the site is to hold the target) and makes a round trip to
check it. Because the robots do not talk to each other, two robots
may independently pick the same site and check it twice. This is
the ``no-coordination floor''---the baseline against which every
smarter model is measured.

In Model 1 each of the $R$ robots independently samples an unvisited
site from the normalized prior and makes a base-to-site round trip of
cost $2\,d(x_0^r, x_i)$. There is no communication, so robots may
select the same site in the same round. The purpose of Model 1 is to
quantify the cost of no coordination at all and to serve as a floor
against which every other model is measured.

\subsection{Model 2: Unconstrained-Energy Auction}

\paragraph{The idea.} Model 2 is the other extreme from Model 1:
robots share full information and pretend battery is not a problem.
The auction is the coordination mechanism. Each robot advertises a
bid for every unvisited site, reflecting how attractive that site is
from its current position (a nearby high-probability site gets a
high bid; a far low-probability site gets a low one). The auction
matches robots to sites greedily by bid, so two robots never
accidentally go to the same place. Movement is ``straight line from
where I am now to the next site''---no returning to base between
visits---because the battery is imagined to be infinite. Model 2
therefore tells us what is possible when the only constraint is
coordination. The gap between M1 and M2 is the cost of not
coordinating. The gap between M2 and the battery-constrained models
that come next (M3, M4, M4*) is the cost of the battery.

In Model 2 the robots share full information and execute a sequential
single-item (SSI) auction. Each robot $r$ computes a bid for every
unvisited site $i$ using its current position $\mathrm{pos}_r$
(initially $x_0^r$):
\[
  b_{r,i} = \frac{p_i}{d(\mathrm{pos}_r, x_i)}.
\]
The pair $(r^*, i^*) = \arg\max_{r,i} b_{r,i}$ wins the first auction
slot; both $r^*$ and $i^*$ are then removed and the auction repeats
until all robots have assignments. Winners travel one-way to their
assigned sites---there is no base return and no energy limit---and
their positions update for the next round. Model 2 is a strong
baseline but \emph{not} an algorithmic lower bound, because the greedy
SSI order can accept chains that a smarter global planner would
reject. In the energy sweep of Chapter~V the multi-sortie planner
Model 4* actually overtakes Model 2 for large $E$, confirming that
Model 2 is not an omniscient best.

\subsection{Model 3: Auction with Single-Node Sorties}

\paragraph{The idea.} Model 3 turns the battery back on. Now each
robot has a fixed energy budget $E$ and must return to base after
each site visit. The auction runs at base at the start of every
round: robots share what they have learned so far, bid on
energy-feasible sites, and get assigned one site each. Then they
fly to their assigned site, check it, and fly back to recharge.
Model 3 exposes what this thesis calls the \emph{energy penalty}. A
robot with energy $E = 14$ that is assigned a site at distance $5$
spends $10$ units on the flight there and back, wasting the
remaining $4$ units of battery on nothing useful. Those $4$ units
of waste, summed over many sorties across many rounds, are
substantial. Model 4 is designed to reclaim them.

In Model 3 the SSI auction runs at base each round, restricted to
energy-feasible candidates satisfying $2\,d(x_0^r, x_i) \leq E$. Each
robot executes one round-trip sortie per iteration. This model
exposes the \emph{energy penalty}: a single sortie spends most of its
budget on the two base legs, leaving little for productive search. A
robot with $E = 14$ and an assigned site at distance $5$ spends $10$
units on travel and wastes the remaining $4$. The next model addresses
exactly this waste.

\subsection{Model 4 and Model 4*: Multi-Node Sorties}\label{sec:foreclosure}

\paragraph{The idea.} The fix for Model 3's wasted battery is to let
a robot chain several sites into one flight before returning to
base. The auction still runs at base to pick each robot's \emph{first}
site for the round---this is called the \emph{anchor} site. After
that, while the robot has battery remaining, it picks its next site
greedily: of the unvisited sites still reachable given the current
position and the energy that would be needed to return to base, go
to the one with the highest bid. Chain until no more sites are
reachable, then fly home. The result is much better battery
utilization---a sortie under Model 4 typically visits five or six
sites instead of one.

Model 4* is Model 4 with one single change: the bid function. Where
Model 4 uses the classical $p_i / d$ bid (``rank sites by
probability per unit distance''), Model 4* uses $p_i / d^2$
(``rank sites by probability per unit distance-squared''). The
difference is subtle on paper but significant in effect. Under the
quadratic rule, a site twice as far away looks four times less
attractive instead of twice less, which makes the robot prefer
short chains over long ones. The intuition for why this matters is
the \emph{cost-foreclosure argument} expanded below.

In Model 4, after the auction assigns an initial (anchor) site to
each robot, a \emph{greedy chain extension} appends additional sites
to each robot's tour while the remaining energy permits. Let
$\mathrm{pos}_r$ be the robot's current position (updated after each
hop) and $\mathcal{C}$ the set of sites already claimed by any robot's
chain. The next extension site for robot $r$ is
\begin{equation}\label{eq:chain}
  j^* = \argmax_{j \notin \mathcal{C}} \; \frac{p_j}{d(\mathrm{pos}_r, x_j)},
  \quad \text{subject to} \quad d(\mathrm{pos}_r, x_j) + d(x_j, x_0^r) \leq E^r_{\mathrm{rem}},
\end{equation}
where $E^r_{\mathrm{rem}}$ is the robot's remaining energy. Extensions
proceed in round-robin order across robots until no more sites can be
appended. Model 4* uses the same structure as Model 4 but replaces
$p/d$ with $p/d^2$ in both the initial auction and the chain extension.

The cost-foreclosure argument for the quadratic penalty is as follows.
A sortie that anchors at distance $d$ consumes $2d$ of the budget on
its base legs alone, leaving $E - 2d$ for chain hops. The number of
chain hops possible is approximately $(E - 2d) / d_{\mathrm{hop}}$,
which is a linear function of $d$ with negative slope. The
\emph{product} of (anchor probability) and (expected total progress of
the sortie) therefore favors a superlinear penalty on $d$, which the
bid $p/d^2$ approximates.

Figure~\ref{fig:m4star-walkthrough} shows one full iteration of Model
4* on a small instance, and Algorithm~\ref{alg:m4} gives pseudocode
for the same iteration. The outer loop repeats until the target is
found.

\begin{figure}[ht]
  \centering
  \includegraphics[width=1.0\textwidth]{fig_m4star_walkthrough.png}
  \caption{Model 4* executing one iteration on a ten-site, three-robot instance. (A) Initial priors; site size is proportional to probability. (B) The SSI auction assigns each robot a first (anchor) site. (C) Greedy chain extensions fill the remaining energy budget. (D) After the round, visited sites are removed and the Bayesian posterior concentrates on the unvisited sites.}\label{fig:m4star-walkthrough}
\end{figure}

\begin{algorithm}[ht]
\caption{Model 4 / 4* iteration.}\label{alg:m4}
\begin{algorithmic}[1]
\Require sites $\mathcal{N}$, prior $\mathbf{p}$, bases $\mathcal{B}$, energy $E$, bid $\phi(p, d)$
\State $\mathcal{A} \gets \{i \in \mathcal{N} : \exists r, 2\,d(x_0^r, x_i) \leq E\}$
\State \textbf{auction:} for each $r$, solve $(r, i^r) \gets \argmax \phi(p_i, d(x_0^r, x_i))$ via SSI over $\mathcal{A}$
\State $\mathcal{C} \gets \{i^r : r = 1,\ldots,R\}$
\For{each robot $r$}
  \State $\mathrm{pos}_r \gets x_{i^r}$; $E^r_{\mathrm{rem}} \gets E - d(x_0^r, x_{i^r})$
\EndFor
\While{any robot can extend its chain}
  \For{each robot $r$ in round-robin order}
    \State $j^* \gets \argmax_{j \notin \mathcal{C}} \phi(p_j, d(\mathrm{pos}_r, x_j))$ s.t.\ $d(\mathrm{pos}_r, x_j) + d(x_j, x_0^r) \leq E^r_{\mathrm{rem}}$
    \If{$j^*$ exists}
      \State append $j^*$ to robot $r$'s tour; $\mathcal{C} \gets \mathcal{C} \cup \{j^*\}$
      \State $\mathrm{pos}_r \gets x_{j^*}$; $E^r_{\mathrm{rem}} \gets E^r_{\mathrm{rem}} - d(\mathrm{pos}_r, x_{j^*})$
    \EndIf
  \EndFor
\EndWhile
\State execute all tours; if any robot finds $\tau$ return success
\State Bayesian update on $\mathbf{p}$ using visited sites; repeat
\end{algorithmic}
\end{algorithm}

\subsection{Model 4* Implementation}

The following pseudocode shows one complete auction iteration of Model~4*.
Reading through this code is the fastest way to understand exactly how the
bid function and chain extension work together. Annotations explain each
step.

\begin{verbatim}
def m4star_iteration(sites, p, bases, R, E, visited):
    """
    One auction iteration of Model 4* (p/d^2 bid, greedy chain).
    Input:  sites    - array of n site positions
            p        - current prior probabilities (updated each round)
            bases    - array of R robot base positions
            R        - number of robots
            E        - energy budget per robot per sortie
            visited  - set of already-inspected site indices
    Returns: tours   - list of R site-index lists (one tour per robot)
    """
    unvisited = [i for i in range(len(sites)) if i not in visited]
    pos    = [bases[r].copy() for r in range(R)]  # current position = base
    E_rem  = [E] * R                              # remaining energy per robot
    claimed = set()                               # sites already assigned
    tours  = [[] for _ in range(R)]              # output: one tour per robot

    # --- Phase 1: SSI auction assigns each robot one anchor site ---
    # The robot with the highest p[i]/d^2 bid wins site i.
    # This is run sequentially: once a robot wins, it is removed from bidding.
    available = [i for i in unvisited if any(
        2 * dist(bases[r], sites[i]) <= E for r in range(R))]
    for _ in range(R):                            # at most R winners
        best_bid, best_r, best_i = -inf, None, None
        for r in range(R):
            if tours[r]:                          # already has an anchor
                continue
            for i in available:
                if i in claimed:
                    continue
                d = dist(bases[r], sites[i])
                if 2 * d <= E:                    # round-trip feasibility
                    bid = p[i] / d**2             # the p/d^2 bid function
                    if bid > best_bid:
                        best_bid, best_r, best_i = bid, r, i
        if best_r is None:
            break                                 # no more feasible sites
        tours[best_r].append(best_i)
        claimed.add(best_i)
        pos[best_r]   = sites[best_i]            # robot moves to anchor
        E_rem[best_r] -= dist(bases[best_r], sites[best_i])

    # --- Phase 2: Greedy chain extension ---
    # Each robot greedily adds the next site with the best p/d^2 score,
    # as long as it can still return to base afterward.
    changed = True
    while changed:
        changed = False
        for r in range(R):
            if not tours[r]:
                continue
            best_bid, best_j = -inf, None
            for j in unvisited:
                if j in claimed:
                    continue
                d_next = dist(pos[r], sites[j])       # cost to reach j
                d_home = dist(sites[j], bases[r])     # cost to return home
                if d_next + d_home <= E_rem[r]:       # energy check
                    bid = p[j] / d_next**2
                    if bid > best_bid:
                        best_bid, best_j = bid, j
            if best_j is not None:
                tours[r].append(best_j)
                claimed.add(best_j)
                E_rem[r] -= dist(pos[r], sites[best_j])
                pos[r]    = sites[best_j]
                changed   = True                      # keep extending
    return tours
\end{verbatim}

After each auction iteration, the outer loop performs a Bayesian update:
every inspected site has its prior set to zero and the remaining priors
are renormalized. The loop repeats until the target site appears in a
robot's tour. The full simulation code, including all five models and
all eight experiments, is available at
\url{https://github.com/foojanbabaeeian/Multi-Robot-Algo}.

\subsection{Hungarian Baselines}

The Hungarian algorithm~\citep{Kuhn1955} computes the globally
minimum-cost assignment of robots to tasks in $O(n^3)$. This thesis
evaluates two Hungarian variants as centralized baselines.

The first variant, H$_d$, minimizes total base-to-site distance with
edge cost $c_{r,i} = d(x_0^r, x_i)$. It \emph{ignores} the prior.
This tests whether distance-optimal assignment alone suffices.

The second variant, H$_{p/d}$, gives the Hungarian solver the same
probability information that the distributed auction uses, with edge
cost $c_{r,i} = -p_i / d(x_0^r, x_i)$ (negated to convert the
maximum-bid problem to a minimum-cost problem). Comparing M3 to H$_d$
and H$_{p/d}$ jointly disentangles two previously conflated effects:
the role of probability-awareness in the bid function versus the role
of centralization in the allocation mechanism.

\section{Theoretical Analysis}

This section derives upper bounds on the expected FCR of each model
in closed form, expressed in terms of the instance parameters
$(n, R, E, L)$ and the geometric constants $d_{\mathrm{avg}}$ and
$d_{\mathrm{hop}}$.

\paragraph{Why Lemma~\ref{lem:rounds} is the right starting point.}
Every bound derived in this chapter is built around an estimate of
how many \emph{rounds} the algorithm takes to find the target. A
round is one iteration of the auction-plus-execute loop---so if a
search takes five rounds, the robot team has run the auction five
times. Once we know how many rounds the algorithm takes on
average, converting to total distance traveled by the finder is a
matter of multiplying by the cost of one round (which depends on
the movement pattern of the specific model). Lemma~\ref{lem:rounds}
gives the round count in closed form under coordination---that is,
assuming the $R$ robots check $R$ distinct sites every round, which
is exactly what the auction guarantees.

\begin{lemma}[Expected Rounds under Coordination]\label{lem:rounds}
If $R$ robots check $R$ distinct sites per round, then for any prior
$\mathbf{p}$ the expected number of rounds to find the target
satisfies
\[
  \mathbb{E}[K] \;\leq\; \frac{1}{n}\sum_{i=1}^{n}\left\lceil \frac{i}{R} \right\rceil \;=\; \frac{n + R}{2R} + O(1/n).
\]
\end{lemma}

\begin{proof}
Let $\sigma$ be the search order induced by the auction (sites visited
round by round in decreasing bid). The target at $\sigma$-rank $i$ is
found in round $\lceil i/R \rceil$. Therefore the expected number of
rounds is $\sum_{i=1}^n \Pr[\sigma\text{-rank of target} = i] \cdot
\lceil i/R \rceil$, and this sum is maximized by the uniform
distribution $p_i \equiv 1/n$. Summing
$\frac{1}{n}\sum_{i=1}^n \lceil i/R \rceil$ in blocks of $R$ yields
the stated closed form after an $O(1/n)$ correction. When the prior
concentrates on high-probability sites, the auction checks those
first, and the expected rounds are strictly smaller; the bound is
therefore conservative for non-uniform priors.
\end{proof}

\paragraph{What Theorem~\ref{thm:m1} says in words.} For the
uncoordinated baseline M1, the probability that some robot picks the
target in any given round is small (one in $n$ per robot,
independent across robots). The number of rounds before one of them
gets lucky is therefore governed by a geometric distribution, whose
mean is the reciprocal of the per-round success probability. Every
round before the last one costs a round trip (cost $2\,d_\mathrm{avg}$),
and the final round costs a one-way trip to the target. The bound
is the sum of those costs divided by the optimal cost.

\begin{theorem}[Model 1 Bound]\label{thm:m1}
Let $P = 1 - (1 - 1/n)^R$ be the per-round success probability, and
let $\mathbb{E}[K_1] = 1/P$. Then
\[
  \mathbb{E}[\mathrm{FCR}_1] \leq \frac{(2\,\mathbb{E}[K_1] - 1) \cdot d_{\mathrm{avg}}}{d_{\mathrm{opt}}}.
\]
\end{theorem}

\begin{proof}
Each of $R$ robots independently picks the target with probability
$1/n$. The probability that none succeeds in a round is $(1 - 1/n)^R$,
so per-round success probability is $P$ and the number of rounds is
geometric with $\mathbb{E}[K_1] = 1/P$. The finder makes $K_1 - 1$
round trips, each of expected cost $2\,d_{\mathrm{avg}}$, and one
final one-way trip of expected cost $d_{\mathrm{avg}}$, giving
$\mathbb{E}[D_f] \leq (2\,\mathbb{E}[K_1] - 1) \cdot d_{\mathrm{avg}}$.
\end{proof}

\begin{theorem}[Model 2 Bound]\label{thm:m2}
$
  \mathbb{E}[\mathrm{FCR}_2] \leq \mathbb{E}[K] \cdot d_{\mathrm{hop}} / d_{\mathrm{opt}}.
$
\end{theorem}

\begin{proof}
Under coordination, Lemma~\ref{lem:rounds} gives $\mathbb{E}[K]$ as
the expected number of auction rounds. Robots in Model 2 move
node to node (no base return), and we approximate the expected
distance between consecutive auction winners by the characteristic
inter-site spacing $d_{\mathrm{hop}} = L/\sqrt{n}$. The finder
accumulates at most $\mathbb{E}[K]$ such hops.
\end{proof}

\begin{theorem}[Model 3 Bound]\label{thm:m3}
$
  \mathbb{E}[\mathrm{FCR}_3] \leq (2\,\mathbb{E}[K] - 1) \cdot d_{\mathrm{avg}} / d_{\mathrm{opt}}.
$
\end{theorem}

\begin{proof}
Model 3 has the same round-trip structure as Model 1 (base $\to$ site
$\to$ base), but the auction ensures that $R$ distinct sites are
visited per round, so the coordinated $\mathbb{E}[K]$ of
Lemma~\ref{lem:rounds} replaces the geometric $\mathbb{E}[K_1]$ of
Theorem~\ref{thm:m1}.
\end{proof}

\paragraph{What Theorem~\ref{thm:m4} says in words.} In Model 4
each round is a \emph{multi-site sortie} rather than a single-site
trip, so a round costs (at most) the full energy budget $E$ rather
than twice the distance to one site. The bound for the
\emph{number} of sites checked per round is $S$, counting one anchor
site plus some number of chain hops. With $R$ robots each checking
$S$ sites per round, the team covers $RS$ sites per round, and the
expected number of rounds to find the target falls by a factor of
$S$ compared to Model 3. The worst-case total travel is $E$ per
round times the expected number of rounds.

\begin{theorem}[Model 4 Bound]\label{thm:m4}
Define the expected number of sites per sortie as
\[
  S = \max\!\big(1,\; 1 + \lfloor (E - 2\,d_{\mathrm{avg}}) / d_{\mathrm{hop}} \rfloor\big),
\]
reflecting one anchor site at expected base-leg cost $2\,d_{\mathrm{avg}}$
plus as many hop extensions as the remaining budget allows. Then
\[
  \mathbb{E}[\mathrm{FCR}_4] \leq \frac{E \cdot \mathbb{E}[K_4]}{d_{\mathrm{opt}}},
  \quad \mathbb{E}[K_4] \leq \frac{1}{n} \sum_{i=1}^n \left\lceil \frac{i}{RS} \right\rceil.
\]
\end{theorem}

\begin{proof}
A sortie anchored at site $s_1$ spends $d(x_0^r, x_{s_1}) + d(x_{s_k}, x_0^r)$
on the first and last legs, whose expected sum is approximately
$2\,d_{\mathrm{avg}}$. The remaining energy $E - 2\,d_{\mathrm{avg}}$
supports $\lfloor (E - 2\,d_{\mathrm{avg}})/d_{\mathrm{hop}}\rfloor$
hops of cost $d_{\mathrm{hop}}$ each, giving $S$ total sites. With $R$
robots each visiting $S$ sites, $RS$ sites are eliminated per round,
so by Lemma~\ref{lem:rounds} applied with block size $RS$ the
expected number of rounds is at most
$\frac{1}{n}\sum_{i=1}^n \lceil i / (RS) \rceil$. Each sortie costs at
most $E$, so the finder's expected travel is at most
$E \cdot \mathbb{E}[K_4]$.
\end{proof}

\begin{corollary}[Improvement Factor]\label{cor:improvement}
Model 4 reduces the expected number of rounds by a factor of $S$
relative to Model 3 whenever $S$ sites remain reachable per sortie.
As $E \to \infty$, $S \to \infty$ and Model 4's per-round site
throughput approaches that of the unconstrained node-to-node regime
of Model 2.
\end{corollary}

\subsection{Fleet-Size Scaling of the Expected FCR}\label{sec:scaling}

This subsection makes precise why auction coordination reduces FCR by
a factor of $\Theta(\sqrt{n})$ regardless of fleet size---a result
that explains the robot-sweep behavior observed in Chapter~V and
converts the scaling conjecture of the literature into a proved theorem.

\paragraph{Key Lemma: $\mathbb{E}[D_{\mathrm{opt}}] = L/(2\sqrt{R}) + O(L/R)$.}
Let $R$ robot bases be drawn independently and uniformly from $[0,L]^2$,
and let $D_{\mathrm{opt}} = \min_{r=1}^{R} d(x_0^r, \tau)$ be the
oracle distance to a fixed interior target $\tau$.
The expected oracle distance satisfies
\[
  \mathbb{E}[D_{\mathrm{opt}}] = \frac{L}{2\sqrt{R}} + O\!\left(\frac{L}{R}\right),
\]
proved via the survival function representation
$\mathbb{E}[D_{\mathrm{opt}}] = \int_0^\infty (1 - F(t))^R\,dt$,
where $F(t) = \pi t^2/L^2$ for a disk of radius $t$ in the interior.
Substituting $u = t\sqrt{R\pi}/L$ and applying dominated convergence
with $(1-u^2/R)^R \to e^{-u^2}$ yields the leading term $L/(2\sqrt{R})$
and a correction of $O(L/R)$. The scaling $\mathbb{E}[D_{\mathrm{opt}}]
= \Theta(L/\sqrt{R})$ holds uniformly over all $\tau \in [0,L]^2$.

\paragraph{M1 Scaling: $\mathbb{E}[\mathrm{FCR}_1] = \Theta(1/\sqrt{R})$ (experimental regime).}
In the experimental regime $1 \leq R \leq n$ (specifically $R \in [1,6]$,
$n = 30$ throughout this thesis), the per-round success probability is
$P_{\mathrm{succ}} \approx R/n$, giving $\mathbb{E}[K_1] \approx n/R$
and $\mathbb{E}[D_f] \approx 2(n/R)\,d_{\mathrm{avg}}$.
Using the reciprocal moment $\mathbb{E}[1/D_{\mathrm{opt}}] = \pi\sqrt{R}/L$
(computed from the same survival function argument), and independence of
$K_1$ from the base positions:
\[
  \mathbb{E}[\mathrm{FCR}_1]
  \;\approx\; \frac{2n}{R}\cdot d_{\mathrm{avg}} \cdot \frac{\pi\sqrt{R}}{L}
  \;=\; \frac{2\pi n\,d_{\mathrm{avg}}}{L\,\sqrt{R}}
  \;\approx\; \frac{3.28\,n}{\sqrt{R}}.
\]
Thus $\mathbb{E}[\mathrm{FCR}_1] = \Theta(1/\sqrt{R})$ for fixed $n$
with $R \leq n$.

\paragraph{M2 Scaling and Coordination Gain.}
Under Model 2, Lemma~\ref{lem:rounds} gives $\mathbb{E}[K_2] \leq n/(R) + 1/2$,
and node-to-node hops cost $d_{\mathrm{hop}} = L/\sqrt{n}$ each, yielding
$\mathbb{E}[D_f] = O(L\sqrt{n}/R)$.
The same reciprocal-moment calculation gives:
\[
  \mathbb{E}[\mathrm{FCR}_2] \leq \frac{\pi\sqrt{n}}{\sqrt{R}} \cdot C_J,
  \quad C_J = \frac{\pi}{2},
\]
so $\mathbb{E}[\mathrm{FCR}_2] = \Theta(1/\sqrt{R})$---the \emph{same exponent}
as M1 but with constant factor $\approx \pi\sqrt{n}$ versus M1's $\approx 3.28n$.
Taking the ratio:
\[
  \frac{\mathbb{E}[\mathrm{FCR}_2]}{\mathbb{E}[\mathrm{FCR}_1]}
  \;=\; \Theta\!\left(\frac{1}{\sqrt{n}}\right),
\]
independent of $R$. At $n = 30$, this gives
$\approx 1/(2 \times 0.5214 \times \sqrt{30}) \approx 1/5.7$,
consistent with the observed $6.6\times$ coordination gain (the
small gap arises from the Jensen correction $C_J = \pi/2$, which is
absorbed into the $\Theta$ notation).

\paragraph{What this proves.}
The coordination advantage of $\Theta(\sqrt{n}) \approx 5.7\times$ at
$n = 30$ is proved to be \emph{independent of fleet size} for all
$R \leq n$. This rigorously confirms that communication mechanism,
not the number of robots deployed, is the primary design lever in
multi-robot search. The scaling analysis converts the fleet-size
conjecture raised in prior work into a proved result.

\end{doublespacing}
